\subsection{Extreme rays without vertex zeros}
\label{sec:three:edge-classification}

The strict five-contact family can be characterized without enumerating
contact positions. Apart from affine squares, it accounts for every
extreme ray with positive square coefficients and no vertex zero.

Write $q(x)=q_0+a^\top x+x^\top Qx$, so that $Q_{ii}=b_i$ and
$2Q_{ij}=c_{ij}$ in the coefficient convention of
\Cref{sec:setting}.

\begin{theorem}\label{three:edge-classification}
Suppose that $q$ generates an extreme ray of $\mathcal P_3$ and satisfies
\begin{equation}\label{eq:three:edge-hypotheses}
 b_i>0,\qquad
 q(v)>0\quad(v\in\{0,1\}^3).
\end{equation}
Then $q$ is an affine square, or a cube symmetry transforms it into
$q_{h,d,k}$ with
\[
 d_1,d_2,d_3,k>0,\qquad
 0<h<\min(d_1,d_2),\qquad d_1+d_2-h+k<d_3.
\]
Conversely, every symmetry copy of a family member in this parameter range
generates an extreme ray and satisfies \eqref{eq:three:edge-hypotheses}.
An affine square satisfying \eqref{eq:three:edge-hypotheses} generates an
extreme ray if and only if its affine factor takes both signs on the cube.
In the non-square case, every $2\times2$ principal determinant of $Q$ is
strictly negative.
\end{theorem}

\begin{proof}[Proof outline]
For a nonpositive mixed-coefficient product, the attained dual of the
exact submodular SDP decomposes $q$ into affine squares, nonnegative
coordinate products, diagonal caps, and a nonnegative constant. Its
extremality and positive vertex values force it to be an affine square.
In the remaining sign class, complementing a coordinate if necessary
makes all mixed coefficients positive. A non-square extreme ray cannot
have a stationary zero of a
positive-semidefinite facet restriction: all its zeros would lie in an
affine plane, allowing subtraction of a nonzero affine square.
\Cref{three:stationary-facet} proves this fact, including stationary zeros
on facet boundaries. It excludes facet-interior zeros and forces positive
inward derivatives at every edge zero. An extreme ray must therefore have
at least five edge zeros: fewer give at most eight
linear value and tangency equations in the ten quadratic coefficients,
leaving a nontrivial two-sided perturbation.

The edges containing zeros form a subgraph of the cube. Two such edges
meeting at a vertex must point toward the same Hamming-weight level;
three meeting edges would decompose the quadratic into an affine square
and nonnegative products. The remaining graph lies in the two outer
three-edge stars and the central six-edge cycle. A direct classification
of graphs with at least five edges leaves the five-contact family. Five
central-cycle edges instead give an affine square plus a nonnegative
multiple of
\[
 (1-x_1)(1-x_2)(1-x_3)+x_1x_2x_3,
\]
which is in $\mathcal D_3$; every other graph is excluded by two-variable
nonnegativity. The complete perturbation and graph arguments are given in
\Cref{app:three:edge-classification}. The converse follows from the
five-contact extremality proof and its coefficient formulas. The same
appendix proves the stated characterization of the affine-square rays.
\end{proof}

Thus an extreme ray with positive square coefficients and no vertex zero
is already accounted for by the disjoint cone or the family. Under these
hypotheses, if one $2\times2$ principal determinant of $Q$ is nonnegative,
the ray must be an affine square.

By \Cref{three:reduction}, an extreme ray missing from $\mathcal S$ has
positive square coefficients. The extremality argument preceding
\eqref{eq:three:completeness} makes it extreme in $\mathcal P_3$ as well.
The theorem therefore confines the full completeness question to rays
with a vertex zero.
